% =========================================================================
% SEGMENTAÇÃO INTERATIVA DE IMAGENS VIA GRAPH CUTS
% =========================================================================
\chapter{Aplicação Prática: Segmentação Interativa de Imagens via Cortes em Grafos}\label{sec:segmentacao}

O problema de Segmentação Interativa de Imagens em Visão Computacional particiona uma imagem digital em \textit{foreground} (objeto) e \textit{background} (fundo) minimizando uma função de energia que combina coerência com sementes fornecidas pelo usuário e continuidade entre pixels vizinhos~\cite{cook1998combinatorial}.


\section{Modelagem Energética e Formulação de Rede de Fluxo}

Seja $I$ uma imagem com $n = W \times H$ pixels. A rede de segmentação $D = (V, A)$ é construída como segue:

\begin{enumerate}
    \item \textbf{Vértices:} Cada pixel $p_i$ da imagem corresponde a um vértice, com adição da superfonte $s$ (\textit{foreground}) e do supersumidouro $t$ (\textit{background}), totalizando $|V| = n + 2$.

    \item \textbf{N-links (Neighborhood Links):} Para cada par de pixels adjacentes $(p, q)$ sob 4-conectividade, cria-se um arco bidirecional com peso proporcional à similaridade de cor:
    \begin{equation}
        w_{n}(p, q) = K \cdot \exp\left(-\frac{\|I_p - I_q\|^2}{2\sigma^2}\right)
    \end{equation}
    onde $\|I_p - I_q\|$ é a distância euclidiana no espaço RGB, $\sigma$ é o desvio padrão de sensibilidade e $K$ é a constante de escala. Pixels similares recebem peso elevado; descontinuidades de cor recebem peso reduzido.

    \item \textbf{T-links (Terminal Links):} Para cada pixel $p$, inserem-se os arcos $(s, p)$ com capacidade $w_s(p)$ e $(p, t)$ com capacidade $w_t(p)$.

    Para as sementes fornecidas pelo usuário:
    \begin{itemize}
        \item Se $p \in \mathcal{S}_{\text{fg}}$ (semente de \textit{foreground}): $w_s(p) = \infty$ e $w_t(p) = 0$.
        \item Se $p \in \mathcal{S}_{\text{bg}}$ (semente de \textit{background}): $w_s(p) = 0$ e $w_t(p) = \infty$.
    \end{itemize}

    Para os demais pixels, as capacidades são definidas pela distância mínima de cor em relação às sementes:
    \begin{equation}
        w_s(p) = \max\left( K \cdot \exp\left(-\frac{\min_{s_j \in \mathcal{S}_{\text{fg}}} \|I_p - I_{s_j}\|^2}{2\sigma^2}\right), 1 \right)
    \end{equation}
    \begin{equation}
        w_t(p) = \max\left( K \cdot \exp\left(-\frac{\min_{s_j \in \mathcal{S}_{\text{bg}}} \|I_p - I_{s_j}\|^2}{2\sigma^2}\right), 1 \right)
    \end{equation}
    onde o limite inferior unitário evita capacidades nulas.
\end{enumerate}

\begin{figure}[!ht]
\centering
\begin{tikzpicture}[
    x={(-1.2cm, -0.6cm)},
    y={(1.2cm, -0.6cm)},
    z={(0cm, 1.5cm)}
]

    \node[sink node] (t) at (2,2,-3.5) {$t$};

    \foreach \x in {1,2,3} {
        \foreach \y in {1,2,3} {
            \coordinate (cp\x\y) at (\x,\y,0);
        }
    }

    \foreach \x in {1,2,3} {
        \foreach \y in {1,2,3} {
            \draw[->, red!70!black, dashed, line width=0.8pt, draw opacity=0.65] (cp\x\y) -- (t);
        }
    }

    \draw[fill=gray!15, draw=gray!60, fill opacity=0.85, rounded corners=6pt]
        (0.3,0.3,0) -- (3.7,0.3,0) -- (3.7,3.7,0) -- (0.3,3.7,0) -- cycle;

    \begin{scope}[line width=1.1pt, black!75]
      \foreach \x in {1,2,3} {
        \draw[<->] (\x,1,0) -- (\x,2,0);
        \draw[<->] (\x,2,0) -- (\x,3,0);
      }
      \foreach \y in {1,2,3} {
        \draw[<->] (1,\y,0) -- (2,\y,0);
        \draw[<->] (2,\y,0) -- (3,\y,0);
      }
    \end{scope}

    \foreach \x in {1,2,3} {
        \foreach \y in {1,2,3} {
            \node[pixel node] (p\x\y) at (\x,\y,0) {$p_{\x\y}$};
        }
    }

    \node[source node] (s) at (2,2,3.5) {$s$};

    \foreach \x in {1,2,3} {
        \foreach \y in {1,2,3} {
            \draw[->, green!70!black, dashed, line width=0.8pt, draw opacity=0.65] (s) -- (p\x\y);
        }
    }

    \begin{scope}[x={(1cm,0cm)}, y={(0cm,1cm)}, z={(0cm,0cm)}]
        \node[right, font=\footnotesize, green!60!black] at (3.5, 0.7) {T-links $s \to p$};
        \draw[->, green!70!black, dashed, line width=0.8pt] (2.3, 0.7) -- (3.3, 0.7);

        \node[right, font=\footnotesize, text=black] at (3.5, -0.1) {N-links (vizinhança)};
        \draw[<->, line width=1.1pt, black!75] (2.3, -0.1) -- (3.3, -0.1);

        \node[right, font=\footnotesize, red!60!black] at (3.5, -0.9) {T-links $p \to t$};
        \draw[->, red!70!black, dashed, line width=0.8pt] (2.3, -0.9) -- (3.3, -0.9);
    \end{scope}

\end{tikzpicture}

\caption{Modelagem de grade $3 \times 3$ de pixels como rede de fluxo para segmentação: N-links bidirecionais entre vizinhos planares sobre o grid e T-links verticais conectando cada pixel aos terminais $s$ e $t$.}
\label{fig:segmentacao_grafo}
\end{figure}


\section{Segmentação via Corte Mínimo e Otimização Global}

O corte mínimo $(S, T)$ na rede induz a partição ótima dos pixels: $S$ corresponde ao \textit{foreground} e $T = V \setminus S$ ao \textit{background}.

A capacidade do corte reflete dois custos:
\begin{itemize}
    \item \textbf{Penalidade Regional (T-links):} Secciona-se $s \to p$ com custo $w_s(p)$ se $p \in T$; simetricamente, secciona-se $p \to t$ com custo $w_t(p)$ se $p \in S$.
    \item \textbf{Penalidade de Fronteira (N-links):} Secciona-se o arco entre $p$ e $q$ com custo $w_n(p, q)$ sempre que $p \in S$ e $q \in T$, penalizando descontinuidades espaciais.
\end{itemize}

O corte mínimo minimiza a soma da energia regional e de fronteira.

\begin{figure}[!ht]
\centering
\begin{tikzpicture}[thick]
  \draw[cut S, rounded corners=12pt] (-4.2, 1.6) rectangle (8.2, 6.6);
  \node[cut label S, anchor=west] at (-3.6, 6.0) {$\mathbf{S}$ \textnormal{\footnotesize(Foreground / Objeto)}};

  \draw[cut T, rounded corners=12pt] (-4.2, -3.6) rectangle (8.2, 1.2);
  \node[cut label T, anchor=west] at (-3.6, -2.8) {$\mathbf{T}$ \textnormal{\footnotesize(Background / Fundo)}};

  \node[source node] (s) at (2, 5.8) {$s$};
  \node[sink node] (t) at (2, -2.8) {$t$};

  \node[flow node, fill=blue!10] (p1) at (0, 3) {$p_1$};
  \node[flow node, fill=blue!30] (p2) at (4, 3) {$p_2$};
  \node[flow node, fill=orange!20] (p4) at (0, 0) {$p_4$};
  \node[flow node, fill=orange!40] (p3) at (4, 0) {$p_3$};

  \draw[-, line width=1.3pt] (p1) -- node[above, font=\footnotesize, fill=white, inner sep=1.2pt]{$w_n$ alto} (p2);
  \draw[-, line width=1.3pt] (p4) -- node[below, font=\footnotesize, fill=white, inner sep=1.2pt]{$w_n$ alto} (p3);

  \draw[-, line width=1.3pt] (p1) -- node[edge lbl, right=3pt, pos=0.25]{$w_n$ baixo} (p4);
  \draw[-, line width=1.3pt] (p2) -- node[edge lbl, left=3pt, pos=0.25]{$w_n$ baixo} (p3);

  \draw[->,green!60!black,dashed] (s) -- node[above left, edge lbl green]{$w_s=\infty$} (p1);
  \draw[->,green!60!black,dashed] (s) -- node[above right, edge lbl green]{$w_s$ alto} (p2);

  \draw[->,green!60!black,dashed] (s) to[out=180, in=135] node[left, edge lbl green, pos=0.5]{$w_s$ baixo} (p4);
  \draw[->,green!60!black,dashed] (s) to[out=0, in=45] node[right, edge lbl green, pos=0.5]{$w_s$ baixo} (p3);

  \draw[->,red!60!black,dashed] (p4) -- node[below left, edge lbl red]{$w_t=\infty$} (t);
  \draw[->,red!60!black,dashed] (p3) -- node[below right, edge lbl red]{$w_t$ alto} (t);

  \draw[->,red!60!black,dashed] (p1) to[out=225, in=180] node[left, edge lbl red, pos=0.5]{$w_t$ baixo} (t);
  \draw[->,red!60!black,dashed] (p2) to[out=315, in=0] node[right, edge lbl red, pos=0.5]{$w_t$ baixo} (t);

  \draw[cutviolet, line width=1.8pt, dashed] (-4.2, 1.4) -- (8.2, 1.4)
      node[right, font=\footnotesize\bfseries, text=cutviolet, fill=white, inner sep=2.5pt, yshift=3mm] {Corte Mínimo $(S, T)$};
\end{tikzpicture}

\caption{Estrutura do corte mínimo na segmentação: os pixels superiores ($p_1, p_2$) são atribuídos a $S$ e os inferiores ($p_3, p_4$) a $T$. O corte secciona os N-links de menor capacidade e os T-links cruzados, preservando os T-links infinitos das sementes.}
\label{fig:segmentacao_corte}
\end{figure}


\section{Implementação Computacional}

Com $n = W \times H$ pixels, a rede gerada possui $n + 2$ vértices e aproximadamente $6n$ arcos direcionados. A aplicação utiliza o motor de fluxo \lstinline{PushRelabelImproved}, com política \textit{Highest Label} e \textit{Gap Heuristic} ($\mathcal{O}(|V|^2 \sqrt{|A|})$), apresentando baixo tempo de execução em redes em grade~\cite{goldberg1988new}. O código-fonte da aplicação encontra-se documentado no Apêndice~\ref{ap:github}.

